\documentclass[11pt]{article}
\usepackage[margin=0.5in]{geometry}
\usepackage{amsmath}
\usepackage{booktabs}
\usepackage{graphicx}
\usepackage{hyperref}

\title{Time Series Forecasting Project Expose}
\author{Group ID: 140 \\ Members: Rosela Berberi, Anja Koroveshi}
\date{}

\begin{document}
\maketitle

\section{Idea}
The goal of this project is to investigate deep learning approaches for multivariate
time forecasting by predicting future data load and evaluating it on the provided
benchmark dataset. In out approach we ait to compare recurrent and attention-based 
deep learning architectures and measure their feasibility in terms of short-term 
dynamics and long-range dependency capturing.

\section{Planned architectures}
We are considering two main model families for implementation and evaluation:

1. LST + attention

The first approach is an encoder-decoder architecture model based on the Long Short-
Term Memory networks combined with an attention mechanism. The motivation behind this
choice comes as a result of LSTM models capturing sequential temporal dynamics on 
different states, while attention modules help focus on the most important timesteps 
of a sequence and generating repeated patterns.

2. Transformer / TCN ....

\section{Additional Dataset}
To check if the model generalizes to other data domains beyond the benchmark, we will
use a multivariate weather forecasting dataset such as the Jena Climate Dataset that 
contains measurements across many factors (temperature, humidity, wind speed, etc). 
The format and the cross-variable dependencies that affect weather make the comparison 
between LSTM-attention and ...... more realistic and effective regarding future predictions. 

\section{Planned improvements}
Potential improvements and experiments on the two architectures might include: teacher 
forcing for autoregression, different dropout and regularization strategies, ensemble 
forecasting, sequence series normalization or learning rate scheduling based on gradient 
behaviour. 

To obtain results, the models will be compared against the provided benchmark baselines mentioned in 
the source code, including: Naive last-value forecast, Lag-24/-168 repeat and the 
Seasonal mean baseline. Additionally performance will be measured using the leader board metrics:
MAE, RMSE, MAPE, etc.

\section{Risks and Open questions}
A challenge that might arise during implementation is the risk of error propagation and accumulation 
across the network as well as the difficulty of forecasting highly noisy data with little correlation.
A general problem is related to the trade-off between model complexity and generalization, running to 
the risk of overfitting due to the dataset size. For all challenges, the main question remains whether
transformer or recurrent architectures perform better under the following circumstances.

\section{Sources}


\end{document}
